% Part: first-order-logic
% Chapter: model-theory
% Section: substructures

\documentclass[../../../include/open-logic-section]{subfiles}

\begin{document}

\olfileid[pt]{mod}{bas}{sub}
\olsection{Sub\printtoken{p}{structure}}

O !!{domain} de !!a{structure}~$\Struct{M}$ pode ser um subconjunto do
domínio de outra~$\Struct{M'}$.  Mas, obviamente, só devemos considerar
$\Struct{M}$ uma ``parte'' de $\Struct{M'}$ se não apenas $\Domain{M}
\subseteq \Domain{M'}$, mas também $\Struct{M}$ e $\Struct{M'}$
``concordarem'' quanto à maneira como interpretam os símbolos da
linguagem, ao menos na parte compartilhada~$\Domain{M}$.

\begin{defn}
\ollabel{defn:substructure}
Dadas !!{structure}s $\Struct M$ e $\Struct M'$ para a mesma
linguagem~$\Lang L$, dizemos que $\Struct M$ é uma
\emph{sub!!{structure}} de $\Struct M'$ e que $\Struct M'$ é uma
\emph{extensão} de $\Struct M$, escrito $\Struct M \substruct \Struct M'$,
se, e somente se,
\begin{enumerate}
\item $\Domain{M} \subseteq \Domain{M'}$,
\item Para cada constante $c \in \Lang L$, $\Assign{c}{M} =
    \Assign{c}{M'}$;
\item Para cada !!{function} $n$-ária $f \in \Lang L$,
  $\Assign{f}{M}(a_1, \dots, a_n) = \Assign{f}{M'}(a_1, \dots, a_n)$
  para todos $a_1$, \dots, $a_n \in \Domain{M}$.
\item Para cada !!{predicate} $n$-ário $R \in \Lang L$, $\langle
  a_1, \dots, a_n\rangle \in \Assign{R}{M}$ sse $\langle a_1, \dots,
  a_n\rangle \in \Assign{R}{M'}$ para todos $a_1$, \dots, $a_n \in
  \Domain{M}$.
\end{enumerate}
\end{defn}

\begin{rem}
\ollabel{rem:substructure}
Se a linguagem não contém constantes nem !!{function}s, então qualquer
$N \subseteq \Domain{M}$ determina uma sub!!{structure}~$\Struct{N}$ de
$\Struct M$ com !!{domain}~$\Domain{N} = N$, pondo $\Assign{R}{N} =
\Assign{R}{M} \cap N^n$.
\end{rem}

% provar algo sobre isso? Exemplos?

\end{document}
